\chapter{Distributed serving and state placement}
\label{evod46:chapter}

Where does a live request's state reside? A single process hides that question
behind a pointer. Across hosts, moving the bytes and moving authority become
different operations. Suppose worker a owns state at epoch 7 and worker b
receives a copy. If both accept the next token, the request has forked even
though the copied bytes were perfect. If neither accepts it after a crash,
availability has failed even though confidentiality may remain intact.

Our central contract is that placement may change location but never logical
identity or authority. A serialized cluster model will make preparation,
reservation, verification, commit, cancellation and recovery explicit.
It is a specification for a future distributed implementation, not a network
transport, consensus service, RDMA system or live cluster experiment.

\section{Separate control decisions from data movement}
\index{control plane}\index{data plane}\index{placement}\index{routing}
\label{evod46:control}
The control plane records worker membership, capabilities, capacity, artifacts,
health and topology. The data plane carries token/tensor/state traffic.
Advertisements are evidence with a freshness interval, not permanent facts:
a worker may lose its accelerator after announcing free memory. Workers must
revalidate reservations when accepting work. Authenticate control messages;
being reachable on an internal network is not authority.

A useful advertisement includes device/backend, precision and kernel support,
loaded artifact digests, available memory, zone/rack, schema and health epoch.
Routing chooses a logical implementation; placement chooses a physical
instance of that implementation. The scheduler cannot silently substitute a
different model or precision merely because a worker is idle.

For request $r$, choose from hard-feasible set $\mathcal E(r)$ before ranking:
\[
 w^*=\arg\min_{w\in\mathcal E(r)}
 \bigl(\alpha Q_w+\beta T_{rw}+\gamma C_w+\delta R_w\bigr).
\]
$Q_w,T_{rw}$ are predicted queue/transfer durations in milliseconds;
$C_w,R_w$ are declared cost/risk scores with weights converting to the same
objective units. They are estimates, not measured latency guarantees. Empty
$\mathcal E(r)$ means reject or queue, not relax tenant or artifact constraints.

\Listing{Cluster.eligible}
The fixture requires live membership, exact backend/schema, artifact residency,
tenant permission and an available slot. It then adds a synthetic 10 ms
nonpreferred-zone penalty to queue time and breaks ties by occupancy and ID.
A locality preference is soft; authority and compatibility are hard.

\section{Represent identity independently of location}
\index{ownership epoch}\index{fencing}
\label{evod46:owner}
The handle is $H=(sid,t,d_M,schema,backend,owner,e)$.
Logical ID $sid$, tenant $t$, artifact $d_M$, state interpretation and backend
contract survive migration. Owner and monotonic epoch $e$ change together.
Epoch is a fencing generation, not wall-clock time and not a model-token
position. Value equality with the authoritative directory is required for a
handle to operate in the fixture. This rejects accidental foreign/stale
records; it is not authentication against a hostile process reading the directory.

Immutable weights, tokenizers and compiler outputs can be replicated after
digest verification. Mutable request state instead needs a consistency
protocol. We choose one authoritative owner, with temporary destination copies
that cannot execute. A replicated durable-memory service can have another
consistency model; do not infer it from engine-state ownership.

\Plate{01-owner}{A copied snapshot at b is not an owner. Only the directory transition from epoch 7 to 8 authorizes b and fences a. The identity fields above the timeline remain fixed.}{evod46:ownerplate}

\section{Derive migration from a stable snapshot}
\index{migration}\index{snapshot}\index{reservation}
\label{evod46:migrate}
The package's prepare--copy--verify--commit sequence is retained, but preparing
must also freeze updates or capture a versioned snapshot. Otherwise the source
could advance while the old snapshot is copied and commit a stale continuation.
We choose freezing to make the reference small and exact. Reads remain allowed;
writes are rejected until commit or abort.

\Listing{Cluster.prepare}
Preparation checks the complete destination contract before reserving one
slot. It registers an exact transfer object containing the source handle,
destination, immutable byte snapshot and digest. A production transfer needs
an authenticated envelope binding these fields, size bounds and chunk numbers;
a bare hash over bytes is not a MAC and does not establish who sent them.

Copy verification checks bytes, then commit changes owner and epoch in one
serialized directory transition. Destination capacity was already reserved.
The source reservation is released only at commit; during the transfer both
source storage and destination staging consume capacity.
\Listing{Cluster.commit}

\Plate{02-migrate}{Prepare reserves b and freezes a; a corrupted copy cannot cross the commit point. Abort restores a's writable state and releases b. Successful commit changes directory owner/epoch once and rejects old-epoch operations.}{evod46:migrationplate}

\Needspace{20\baselineskip}
\subsection{Safety proof and its distributed premise}
\begin{proposition}[Single authoritative owner in the serial model]
Starting from one current handle per live state, any sequence of accepted
create, prepare, abort, commit, cancel and recover calls leaves at most one
handle accepted for mutation of each logical state.
\end{proposition}
\begin{proof}
Create installs one handle. Prepare retains it but freezes writes; its copied
payload grants no execution authority. Abort retains the handle and unfreezes
it. Commit replaces it by an otherwise identical handle with a new owner and
epoch. Exact directory comparison rejects the old handle. Cancel marks state
terminal and advances its epoch; recovery replaces a failed owner's handle
with a new epoch. No operation creates a second authoritative directory row.
Induction over calls proves the property.
\end{proof}

This proof assumes one serial, durable-in-concept authority transition and
that every side effect respects its fence. Two partitioned controllers each
claiming authority violate the premise. A distributed implementation needs a
linearizable metadata service, fenced storage operations or another explicitly
proved protocol. Expiring a lease using local clocks does not by itself stop
a paused old worker from resuming and writing. Fencing must be enforced at
the resource that accepts the mutation. Our Python assignment supplies neither
consensus nor crash durability.

\subsection{Conserve both owners and staging reservations}
For worker $w$, let $O_w$ be current noncancelled owners, $T_w$ incoming
prepared transfers and $D_w$ held prefill/decode reservations. With capacity
$C_w$ in fixture slots,
\[
 used_w=|O_w|,\quad reserved_w=|T_w|+|D_w|,
 \qquad used_w+reserved_w\le C_w.
\]
Preparing adds $T_w$; commit exchanges one destination staging slot for an
owner slot and removes the source owner slot. Abort removes only staging.
Cancellation removes the owner and any pending transfer exactly once.
Repeated migrations and fault tests check these equalities against directory
and transfer tables rather than merely checking the final payload.

\section{Failure windows, cancellation and recovery}
\index{cancellation}\index{replay}\index{exactly-once effects}
\label{evod46:failure}
Before prepare, failure leaves no destination obligation. During copy,
destination failure permits abort if the source remains usable; source failure
requires evidence-backed recovery or request failure. After commit, the old
handle is stale even if the acknowledgement was lost. Retrying must query the
authoritative directory by migration identity; blindly replaying commit can
allocate twice. The reference rejects settled transfer objects rather than
implementing a durable retry journal.

Cancellation targets the current owner epoch, aborts a prepared transfer and
makes state terminal. A late old transfer cannot resurrect it. Real
cancellation also requires propagation, device acknowledgement and delayed
physical reclamation; changing a directory flag cannot stop a kernel already
accessing memory. We model only serialized logical state and capacity.

\Needspace{23\baselineskip}
\Listing{Cluster.recover}
Recovery requires a failed active owner, a compatible destination and replay
bytes matching a digest supplied by a trusted checkpoint/replay ledger.
The expected digest must be independent evidence: hashing arbitrary replay
and passing that same hash proves nothing. The fixture checks equality but
implements neither that durable ledger nor deterministic replay itself.
If no valid checkpoint or replay source exists, fail or restart the request;
never guess missing state. An owner failing during a frozen transfer requires
explicit abort before this recovery path.

Internal state replay does not guarantee exactly-once external effects. A tool
may commit an action and crash before recording its acknowledgement. Retrying
can duplicate it unless the external service supports an idempotency or
transaction protocol. The logical model performs no external action.

\section{DOGMA and Hermon move different representations}
\index{DOGMA!state movement}\index{Hermon DNA!KV movement}
\label{evod46:representations}
Chapter 31's restricted diagonal recurrence carries a hidden vector and token
position; it is not full DOGMA. For a width $D$ float-$b$ vector, the hidden
payload is $Db$ bytes, excluding metadata, banks and parameters. A broader
native model may carry much more. Compactness alone says nothing about equal
task capacity or speed.

Chapter 35's causal Hermon adapter uses the actual Chapter 23 multi-layer
Transformer with learned absolute positions. Its K/V payload is
\[
 M_{KV}=2LtH_{kv}d_hb\quad\text{bytes},
\]
for $L$ layers, $t$ consumed tokens, $H_{kv}$ KV heads and width $d_h$.
Migration must preserve ordered per-layer content, absolute consumed position,
owner and artifact/mode metadata. A paged representation adds page tables,
partial tails and references. Paging solves local allocation, not transport.

The local tests serialize actual Chapter 31 hidden values and actual Chapter
35 per-layer K/V into bounded data-only fixture envelopes, transfer them and
reconstruct through those same model interfaces. One-token suffix logits
match the unmoved state exactly. This is a trusted test-local adapter; it is
not an untrusted checkpoint loader, page migration, pool ownership transfer
or an integrated distributed engine. No detached toy replaces either model.

\section{Place prefill and decode together only when justified}
\index{prefill/decode disaggregation}\index{transfer model}
\label{evod46:pd}
Transformer prefill and decode can interfere and favor different parallelism.
DistServe studies their separation, phase-specific allocation and bandwidth-aware
placement~\cite{distserve46}. Its evaluated system is not this oracle, and its
reported improvements cannot be transferred to Hermon or DOGMA without measurement.

Disaggregation introduces state transfer between phases. With payload $B$
bytes, sustained bandwidth $W$ bytes/s and fixed serialization/coordination
cost $\ell$ ms, a simple estimate is
\[
 T_{move}=\ell+1000 B/W\quad\text{ms}.
\]
It excludes contention, packet retries, overlap and reconstruction costs.
Separate when predicted queue/interference savings exceed added critical-path
transfer and coordination under the stated SLO. This is a plan-selection
hypothesis, not a universal performance claim.

\begin{center}\input{\Generated/transfer.tex}\end{center}
The assigned comparison uses 1 GB/s decimal bandwidth and 0.1 ms fixed cost.
The 8,192-byte vector is a width-1,024 float64 payload example. The 384 MiB KV
example uses $L=24,t=4096,H_{kv}=8,d_h=128,b=2$ from Chapter 35's arithmetic
example. These are different representations, not equal-capability models;
neither row measures actual network or device performance.

\Plate{03-pd}{Admission reserves both pools before starting prefill. The KV transfer is an explicit causal dependency before decode; a full decode pool rejects the reservation without stranding prefill work.}{evod46:pdplate}
\Listing{pd_reserve}

Our paired reservation checks both pools before changing either counter, and
returns a ticket released once. It does not start prefill or convert that ticket
into a live engine request. A concurrent coordinator needs transactional or
leased reservations and crash cleanup. Per-tenant quota also remains an
upstream prerequisite, not an implemented distributed quota ledger.

\section{Topology, backpressure and distributed observations}
\index{topology}\index{backpressure}\index{clock synchronization}
\label{evod46:observe}
GPU-to-GPU, host, rack and zone links differ in latency, bandwidth and failure
domain. A faster empty worker can lose to a nearby loaded one after transfer;
a nearby incompatible worker remains ineligible. Measure serialization,
network, destination reconstruction and commit separately, while recording
overlap. Backpressure must expose downstream reservations and queue pressure
upstream; unbounded queues merely move overload and grow retained state.

Each placement/migration record should bind request-safe correlation ID,
tenant scope, artifact, logical backend, worker, owner epoch, reason and
estimated/actual transfer bytes. Trace prefill, transfer and decode as distinct
causal spans. Across hosts, clock offsets can make naive subtraction invalid;
use local durations and explicit clock synchronization/error bounds where
cross-host timestamps are compared. The reference records sequential event
IDs only and does not measure distributed time.

\section{Fault injection is a protocol test, not a slogan}
\label{evod46:tests}
Tests corrupt transfer bytes, copy issued transfer objects, kill source or
destination during copy, cancel before/after preparation and commit, fill
destination capacity, attempt incompatible placements, repeat migrations and
reject unverified replay. Actual native/KV continuation fixtures test byte
preservation at the model boundary. A destination reservation is not counted
as a second owner, and failed preparation must not alter capacity.

Production validation should delay/duplicate messages, lose commit replies,
partition the metadata service, restart processes with persistent journals,
kill a worker at every protocol boundary and race real device completion.
Write the expected owner, epoch, capacity and output disposition for each
fault before running it. No network fault or remote experiment occurs here.

\section{Twelve worked exercises}
\begin{enumerate}
\item \textbf{Routing or placement?} Choosing Hermon versus a native engine is
routing; choosing worker b for that engine is placement. \emph{Solution:}
preserve the semantic choice while optimizing its physical realization.
\item \textbf{Migrate epoch 7.} Which fields change? \emph{Solution:} owner and
epoch, now 8. Logical ID, tenant, artifact, schema and backend remain fixed.
\item \textbf{Corrupt the copy.} When can authority move? \emph{Solution:}
only after verification. Failure retains source authority and staging until abort.
\item \textbf{Use the old handle.} It still has correct bytes.
\emph{Solution:} reject owner/epoch mismatch; numerical validity is not authority.
\item \textbf{Replicate weights versus request state.} \emph{Solution:}
immutable digest-verified weights have no update race. Mutable replicas require
an explicit consistency/ownership protocol.
\item \textbf{Compute movement.} Transfer 1 MB at 1 GB/s with 0.1 ms fixed cost.
\emph{Solution:} $1000(10^6/10^9)+0.1=1.1$ ms, before contention or reconstruction.
\item \textbf{Replay an acknowledged tool?} \emph{Solution:} model-state replay
does not deduplicate an external effect; use the tool's own idempotency contract.
\item \textbf{Lose state without evidence.} \emph{Solution:} fail/restart the
request. A plausible invented state violates continuation correctness.
\item \textbf{Reserve downstream first.} Decode is full but prefill is empty.
\emph{Solution:} reject the pair before acquiring either reservation; otherwise
prefill can produce state that has no admitted destination.
\item \textbf{Backpressure or larger queues?} \emph{Solution:} bound queues and
expose pressure. Larger unbounded queues increase memory and tail latency.
\item \textbf{Compact state guarantees cheaper service?} \emph{Solution:} no;
include actual representation, serialization, weights, computation, capability
and topology. The payload table proves none of these end-to-end comparisons.
\item \textbf{Design distributed validation.} \emph{Rubric:} define hard placement
constraints; directory linearization/fencing; stable snapshots; reservation
accounting; crash/cancel/retry outcomes; trusted replay evidence and tool effects;
measure topology and trace provenance. A serial oracle cannot close those gates.
\end{enumerate}

\section{Terms and the measurement boundary}
\Term{Fencing epoch} rejects operations from an obsolete owner.
\Term{Staging reservation} accounts for a copy that cannot execute.
\Term{Linearization point} is the single authority transition assumed here.
Chapter 47 asks which placement and kernel costs are actually measured, and
which remain arithmetic estimates or synthetic traces.
